\chapter{Requisitos no funcionales}
\label{anexo:rnf}

Mientras que el Anexo~\ref{anexo:casos-uso} trata los requisitos funcionales del sistema en forma de casos de uso, este anexo recoge sus \textbf{requisitos no funcionales}, es decir, las propiedades de calidad que el sistema debe satisfacer independientemente de las funciones concretas que ofrece.
En un sistema de navegación volumétrica el rendimiento y el consumo de memoria condicionan la viabilidad del producto, por lo que estos requisitos son de gran importancia.

Cada requisito se deriva de los objetivos del proyecto (Sección~\ref{sec:objetivos}) y se agrupa según las características de calidad habituales del software, como el rendimiento, la fiabilidad, la usabilidad o la mantenibilidad.
Para cada uno se indica la iteración del Capítulo~\ref{chap:desarrollo} en la que se trata.

\medskip

\renewcommand{\arraystretch}{1.3}
\begin{longtable}{@{}p{1.6cm} p{2.5cm} p{6.4cm} p{1.3cm}@{}}
\caption{Requisitos no funcionales del sistema de navegación volumétrica y la iteración del Capítulo~\ref{chap:desarrollo} en la que se abordan.}
\label{tab:rnf} \\
\rowcolor{udcpink!25}
\textbf{ID} & \textbf{Categoría} & \textbf{Requisito} & \textbf{Iter.} \\
\hline
\endfirsthead
\rowcolor{udcpink!25}
\textbf{ID} & \textbf{Categoría} & \textbf{Requisito} & \textbf{Iter.} \\
\hline
\endhead
\textbf{RNF-01} & Eficiencia de recursos & El espacio navegable debe representarse de forma compacta, de modo que abarque grandes volúmenes tridimensionales sin un coste de memoria prohibitivo. & 1 \\
\hline
\textbf{RNF-02} & Rendimiento & La construcción del volumen (\gls{bake}) debe completarse en un tiempo razonable incluso en escenarios con geometría densa. & 1, 5 \\
\hline
\textbf{RNF-03} & Rendimiento & El cálculo de rutas debe resolver numerosas peticiones simultáneas sobre volúmenes densos sin degradar la tasa de fotogramas. & 2, 5 \\
\hline
\textbf{RNF-04} & Capacidad de respuesta & El cálculo de rutas no debe bloquear el hilo principal del juego y debe poder cancelarse en cualquier momento. & 4 \\
\hline
\textbf{RNF-05} & Fiabilidad & La serialización del volumen debe ser reversible y la detección de escenas obsoletas debe ser reproducible entre ejecuciones distintas. & 1 \\
\hline
\textbf{RNF-06} & Fiabilidad & Toda ruta marcada como libre debe conservar la holgura necesaria para el tamaño del agente y evitar de forma robusta el \textit{clipping} con la geometría. & 3 \\
\hline
\textbf{RNF-07} & Usabilidad & La configuración y la depuración del sistema deben ser accesibles para perfiles no técnicos sin necesidad de programar. & 4 \\
\hline
\textbf{RNF-08} & Mantenibilidad & El sistema debe organizarse en módulos desacoplados que faciliten su extensión y mantengan el código de edición y de pruebas fuera del producto final. & 0 \\
\hline
\textbf{RNF-09} & Portabilidad & El sistema debe empaquetarse según los estándares de Unity para permitir su integración inmediata en otros proyectos. & 4 \\
\hline
\end{longtable}
\renewcommand{\arraystretch}{1.0}
